\section{Paying the information increment by thinning}
\label{thinning:section}

The canonical optimizer with one positive common corner has a useful
feature: its cost is linear in the mass of its interior core.  We will
show that the information-rate increment is at most the core's
increment multiplied by that mass.  At full mass the increment is
paid by the one-bit posterior comparison.  This will settle the whole
high-entropy branch of the canonical cost problem.

The scalar inputs already proved here are
\begin{equation}
\label{thinning:scalar-inputs}
 F(0)=0,\qquad F'(I)\ge2,\qquad
 0<F''(I)\le\frac{M}{(L-I)^2},\qquad
 M=\frac{4L^2}{3},\qquad 2>3M.
\end{equation}
Indeed, the slope follows from
Lemma~\ref{scalar:profile-convexity}, the curvature bound from
Lemma~\ref{scalar:scaled-curvature}, and the strict margin from
Lemma~\ref{scalar:entropy-bounds}.  The proof isolates the roles of
these two bounds: the slope supplies a positive term, while the
curvature controls the possible loss over an entropy interval.

\subsection{The common-output path}

Fix an interior ordered pair $a_c>b_c$ with nonnegative center, and
write
\[
 x_c=\frac{a_c+b_c}{2}\ge0,\qquad
 d_c=\frac{a_c-b_c}{2}>0,\qquad
 h_c=\frac{H(a_c)+H(b_c)}2,\qquad
 J_c=H(x_c)-h_c.
\]
For $0<\lambda\le1$, put mass $\lambda$ on this core and mass
$1-\lambda$ on $(1,1)$.  The second atom is a common output: its
coordinates agree, and it contributes neither entropy nor cost.  The
mixture has means $x+d$ and $x-d$, where
\begin{equation}
\label{thinning:path}
 x=1-\lambda(1-x_c),\qquad d=\lambda d_c,
 \qquad e=\lambda h_c.
\end{equation}
For this path define
\begin{equation}
\label{thinning:information}
 \bar h=\frac{H(x+d)+H(x-d)}2,\qquad
 \nu=\bar h-e,\qquad J=H(x)-\bar h,
 \qquad I=\nu+J=H(x)-e.
\end{equation}
The average child information is $\nu$, and $J$ is the entropy gained
by averaging the two child means.  Since
$0<d_c<1-x_c\le1$, the mixture has $x\ge0$ and
$0<d<1-x$.  Entropy concavity gives $\nu\ge0$; strict concavity
gives $J>0$; and $e>0$ gives $I<L$.  Thus every rate argument lies
in $[0,L)$.

These quantities have a channel interpretation.  Let $X$ be a fair
sign and let $B\in\{0,1\}$ be an independent flag with
$\Pr(B=1)=\lambda$.  Given $X,B$, draw a sign $U$ with conditional mean
\[
 \E[U\mid X,B]=(1-B)+B(x_c+d_cX).
\]
Thus $B=0$ selects the common output and $B=1$ selects the core.
The three relevant entropies are
$H(U)=H(x)$, $H(U\mid X)=\bar h$, and
$H(U\mid X,B)=e$.  The chain rule therefore reads
\begin{equation}
\label{thinning:flag-information}
 I=I(U;X,B),\qquad \nu=I(U;B\mid X),\qquad
 J=I(U;X),\qquad I=\nu+J.
\end{equation}
In particular, $\nu$ is the information about the flag after the
input has been revealed.

\begin{theorem}[Common-output payment]
\label{thinning:payment}
For the core and mass in \eqref{thinning:path},
\begin{equation}
\label{thinning:value}
 F(\nu+J)-F(\nu)\le\lambda F(J_c).
\end{equation}
\end{theorem}

We prove this by showing that $[F(\nu+J)-F(\nu)]/\lambda$
increases with $\lambda$. This normalization matches the mixture's
cost $\lambda c(a_c,b_c)$. We first compute its derivative and
identify the entropy-interval estimate that makes it nonnegative.

\subsection{The normalized derivative}

For the state \eqref{thinning:path}, set
\begin{equation}
\label{thinning:window-data}
 \begin{gathered}
 p=1+x,\qquad \sigma=1-x,\qquad b=\log(1+x),\\
 q=e+\psi(x)=L-I,\qquad s=e+L-\bar h=L-\nu,\\
 \kappa=-\frac12\log\left(1-\frac{d^2}{p^2}\right)>0.
 \end{gathered}
\end{equation}
In particular $s-q=J$ and $0<q<s\le L$. Write
\[
 \Delta(\lambda)=F(I)-F(\nu),\qquad
 A=\log\frac2{1+x}=L-b.
\]
For $h_{\mathrm{bin}}(t)=H(1-2t)$, the entropy identity is
\[
 t h_{\mathrm{bin}}'(t)-h_{\mathrm{bin}}(t)=\log(1-t).
\]
Apply it first to
$I=h_{\mathrm{bin}}(\lambda(1-x_c)/2)-\lambda h_c$,
then to the two child expressions with initial probabilities
$(1-a_c)/2$ and $(1-b_c)/2$, and average. The result is
\begin{equation}
\label{thinning:path-derivatives}
 \lambda I'=I-A,\qquad
 \lambda\nu'=\nu-A-\kappa.
\end{equation}
For the second identity, the two logarithms have average
\[
 \frac12\log\frac{(1+x)^2-d^2}{4}=-A-\kappa.
\]

Substitution gives the exact Euler derivative
\begin{equation}
\label{thinning:normalized-derivative}
 \lambda\Delta'-\Delta
 =\int_\nu^I(u-A)F''(u)\,du+F'(\nu)\kappa.
\end{equation}
Indeed, at the chosen value of $\lambda$, the derivative in $u$
of $(u-A)F'(u)-F(u)$ is $(u-A)F''(u)$. The value of $A$ is
held fixed only during this auxiliary integration; both path
derivatives have already been included.

The last term in \eqref{thinning:normalized-derivative} is at least
$2\kappa$. Convexity makes the integrand nonnegative where $u\ge A$.
On the remaining part, the curvature bound
$F''(u)\le M/(L-u)^2$ and the substitution $h=L-u$ give
\begin{equation}
\label{thinning:derivative-loss}
 \lambda\Delta'-\Delta
 \ge2\kappa-M\int_q^s\frac{(h-b)_+}{h^2}\,dh.
\end{equation}
Thus it suffices to prove that this integral is at most $3\kappa$:
the strict margin $2>3M$ will then make
$(\Delta/\lambda)'$ nonnegative. The cutoff $b$ marks exactly
the part of the entropy interval where the curvature term can be
negative. We will first prove an interval estimate for this kernel,
then check its hypotheses using the entropy of the core.

The positive term also has a channel interpretation. For the channel
in \eqref{thinning:flag-information}, the posterior of the fair input
at the positive output is
\[
 \Pr(X=z\mid U=1)=\frac12\left(1+\frac{zd}{p}\right),
 \qquad z\in\{-1,1\}.
\]
Consequently
\[
 \kappa=\frac12\sum_{z=\pm1}
 \log\frac{1/2}{\Pr(X=z\mid U=1)}
 =D_{\mathrm{KL}}\!\left(
 \operatorname{Unif}\{-1,1\}\,\middle\|\,
 \mathcal L(X\mid U=1)\right).
\]
This is the relative entropy from the prior to the positive-output
posterior, with the flag marginalized. It is an auxiliary path
quantity, distinct from the canonical cost $\mathcal K$.

\subsection{An estimate for entropy intervals}

\begin{lemma}[An interval estimate with a logarithmic cutoff]
\label{thinning:interval}
Let $0\le x<1$, and put
\[
 p=1+x,\qquad \sigma=1-x,\qquad
 \gamma=\frac{\sigma}{p},\qquad
 \beta=\frac{\log(1+x)}L.
\]
Suppose
\[
 0\le z<\gamma,\qquad \sigma-pz\le u\le1-z.
\]
For $\beta>0$, define $k_\beta(v)=(v-\beta)_+/v^2$ for $v>0$
and $k_\beta(0)=0$.  For $\beta=0$, set $k_0(v)=1/v$ for $v>0$.
Then
\begin{equation}
\label{thinning:interval-bound}
 \int_u^{u+z}k_\beta(v)\,dv
 \le-\frac32\log(1-\gamma z).
\end{equation}
If $x>0$, the endpoint $z=\gamma$ is also allowed.  At $z=0$,
both sides are zero.
\end{lemma}

\begin{proof}
We first bound the cutoff by a rational function:
\begin{equation}
\label{thinning:cutoff-bound}
 \beta\ge b_*:=\frac{4x}{3+x},\qquad
 1-\beta\le\frac{3\sigma}{3+x}\le\frac32\gamma.
\end{equation}
For $0<x\le1$, let $R(x)=(3+x)\log(1+x)/x$.  Direct
differentiation gives
\[
 x^2R'(x)=N(x):=\frac{x(3+x)}{1+x}-3\log(1+x),
 \qquad N'(x)=-\frac{x(1-x)}{(1+x)^2}\le0.
\]
Since $N(0)=0$, $R$ decreases and $R(x)\ge R(1)=4L$.
This proves the first inequality, including $x=0$ by continuity.
The other two follow from $1-b_*=3\sigma/(3+x)$ and $x\le1$.

Consider an admissible interval $[u,u+z]$ with $z<\gamma$.
The constraints imply
\[
 a:=\frac{2-u}{1+z}\in[1,p].
\]
Keep this $a$ and the cutoff $\beta$ fixed, and consider the affine
family of intervals
\begin{equation}
\label{thinning:affine-intervals}
 u_a(s)=2-a(1+s),\qquad v_a(s)=u_a(s)+s,
 \qquad 0\le s\le z.
\end{equation}
Because $1\le a\le p$, these intervals satisfy
\[
 0<\sigma-ps\le u_a(s)\le v_a(s)\le1.
\]
The initial interval has length zero; the final one is $[u,u+z]$.
The Leibniz rule gives
\[
 \frac{d}{ds}\int_{u_a(s)}^{v_a(s)}k_\beta(v)\,dv
 =D_a(s):=a k_\beta(u_a(s))-(a-1)k_\beta(v_a(s)).
\]
We claim that
\begin{equation}
\label{thinning:interval-derivative}
 D_a(s)\le\frac{3\gamma}{2(1-\gamma s)}.
\end{equation}

For this estimate, abbreviate $u=u_a(s)$ and $v=v_a(s)$.
Above the cutoff,
$k_\beta'(v)=(2\beta-v)/v^3$; thus $k_\beta$ first increases
and then decreases, with the decreasing case $\beta=0$ included.
Its minimum on $[u,1]$ is at an endpoint.  Consequently
\[
 k_\beta(v)\ge\min\{k_\beta(u),1-\beta\}.
\]
If $k_\beta(u)\le1-\beta$, this gives
$D_a\le k_\beta(u)\le1-\beta$, and
\eqref{thinning:interval-derivative} follows from
\eqref{thinning:cutoff-bound}.

Otherwise set $A_*=k_\beta(u)-(1-\beta)>0$.  Then $u>\beta$,
and
\[
 D_a\le aA_*+(1-\beta),\qquad
 Y:=\frac{1-\gamma s}{\gamma}
 =\frac2\sigma-\frac{2-u}{a}>0.
\]
To bound their product, hold $u,\beta,x$ fixed and vary $a$ in
the algebraic expression
\[
 \Theta(a)=\{aA_*+(1-\beta)\}
 \left\{\frac2\sigma-\frac{2-u}{a}\right\}.
\]
It increases with $a$, since
\[
 \Theta'(a)=\frac{2A_*}{\sigma}
 +\frac{(2-u)(1-\beta)}{a^2}>0.
\]
The actual value has $a\le p$, so
\begin{equation}
\label{thinning:transport-bound}
 YD_a\le\Theta(p)
 =\{p k_\beta(u)-x(1-\beta)\}
   \frac{4x+\sigma u}{\sigma p}.
\end{equation}
This last expression decreases with $\beta$: its positive factor
is independent of $\beta$, and the first factor has derivative
$-p/u^2+x<0$.  Since $u>\beta\ge b_*$, replacing $\beta$ by
$b_*$ keeps $u$ above the cutoff and gives an upper bound.
The hard-case condition also gives
\[
 u^2\{k_\beta(u)-(1-\beta)\}
 =(1-u)\{(1-\beta)u-\beta\}>0.
\]
Thus $u<1$ and
$u>\beta/(1-\beta)\ge b_*/(1-b_*)$.
Putting $t=b_*/u$ now yields $0\le t<1/(1+u)$.
Using $x=3b_*/(4-b_*)$, substitution and a common denominator give
\[
 \frac32-\left.\Theta(p)\right|_{\beta=b_*}
 =\frac{\mathcal P(t,u)}{2(1-tu)(2+tu)},
\]
where the numerator satisfies
\begin{equation}
\label{thinning:polynomial}
 \begin{split}
 \mathcal P(t,u):={}&3(2-tu-t^2u^2)\\
 &-\{2(2+tu)(1-t)-3tu^2(1-tu)\}\{1+t(3-u)\}\ge0,
 \qquad 0\le u\le1,\quad 0\le t\le\frac1{1+u}.
 \end{split}
\end{equation}
To prove this polynomial inequality, put
\[
 C_2=12-8u+8u^2-6u^3,\qquad C_1=3u^2-u-8,\qquad
 C_3=u(3-u)(2-3u^2).
\]
Expansion gives
\begin{equation}
\label{thinning:polynomial-decomposition}
 \mathcal P=2+C_1t+C_2t^2+C_3t^3.
\end{equation}
Since
\[
 C_3+u(1+u)=u(1-u)(7+6u-3u^2)\ge0,
 \qquad t(1+u)\le1,
\]
we have $C_3t\ge-u$: this is immediate for $C_3\ge0$, while for
$C_3<0$ it follows from $C_3\ge-u(1+u)$ and $t\le1/(1+u)$.
Hence $\mathcal P\ge2+C_1t+(C_2-u)t^2$.
The exact identity
\[
 8(C_2-u)-C_1^2
 =4+(1-u)\{2(5u-3)^2+u^2(1+9u)+10\}\ge4
\]
therefore gives
\[
 8\mathcal P\ge(4+C_1t)^2+4t^2>0.
\]
This proves \eqref{thinning:polynomial}, including $t=0$.
The denominator in the preceding gap is positive because
$tu=b_*<1$. Thus $YD_a\le3/2$,
which is \eqref{thinning:interval-derivative}.
Integrating from $0$ to $z$ proves \eqref{thinning:interval-bound}.

For $x=0$, admissibility forces $a=1$ and the family is
$[1-s,1]$, whose lower endpoint is positive for $s<1$.  The
integrand $k_0(v)=1/v$ therefore causes no endpoint problem.
For $x>0$, the cutoff is positive and $k_\beta$ is continuous at
zero with the stated extension.  Also $\gamma<1$, so continuity
along the same family extends the bound to $z=\gamma$.
The zero-length case is immediate.
\end{proof}

\subsection{The entropy interval of the channel}

We now return to the quantities in \eqref{thinning:window-data}.
To apply the interval estimate to $[q,s]$, we need a lower bound
on its position and an upper bound on its length $J$. Both follow
from the entropy of the oriented core.

\begin{lemma}[The entropy interval on a thinning path]
\label{thinning:window}
Every state on \eqref{thinning:path} satisfies
\begin{equation}
\label{thinning:window-bound}
 \int_q^s\frac{(h-b)_+}{h^2}\,dh\le3\kappa.
\end{equation}
\end{lemma}

\begin{proof}
The nonnegative center of the core supplies the essential lower
entropy bound.  Write $h_{\mathrm{bin}}(t)=H(1-2t)$, let
$\sigma_c=1-x_c\le1$, and put
\[
 t=\frac{1-d_c/\sigma_c}{2}\in(0,1/2).
\]
The two negative-output probabilities at the core are
$\sigma_c t$ and $\sigma_c(1-t)$.  Concavity and
$h_{\mathrm{bin}}(0)=0$ imply
\[
 h_{\mathrm{bin}}(\sigma_c t)\ge\sigma_c h_{\mathrm{bin}}(t),\qquad
 h_{\mathrm{bin}}(\sigma_c(1-t))
 \ge\sigma_c h_{\mathrm{bin}}(1-t)=\sigma_c h_{\mathrm{bin}}(t).
\]
Averaging and multiplying by $\lambda$, with
$\sigma=\lambda\sigma_c$ and $d/\sigma=d_c/\sigma_c$, gives
\begin{equation}
\label{thinning:entropy-floor}
 e\ge\sigma H(d/\sigma).
\end{equation}
This is a consequence of the oriented common-output path; it is not
an assertion about every pair law with the same means.
The bound $H(r)\ge L(1-r^2)$ in
Lemma~\ref{scalar:entropy-bounds} now yields
\begin{equation}
\label{thinning:quadratic-floor}
 q\ge e\ge L\left(\sigma-\frac{d^2}{\sigma}\right).
\end{equation}

Expanding the binary entropies in $J$ gives the posterior identity
\[
 J=\frac p2\psi(d/p)+\frac\sigma2\psi(d/\sigma).
\]
Using $\psi(r)\le Lr^2$ on both terms, we obtain
\[
 J\le\frac{Ld^2}{\sigma p}=:Lz,\qquad
 0<z<\frac\sigma p.
\]
The upper bound on $z$ follows from $d<\sigma$.  Together with
\eqref{thinning:quadratic-floor}, these estimates say
\[
 q\ge L(\sigma-pz),\qquad q+J=s\le L,
 \qquad L(\sigma-pz)+Lz=L(\sigma-xz)\le L.
\]
We may therefore enlarge $[q,s]$ to an interval of length $Lz$
inside $[L(\sigma-pz),L]$.  More explicitly, choose its lower
endpoint as $q_*:=\min\{q,L-Lz\}$.  Both candidates are at least
$L(\sigma-pz)$, and $q_*\le q$; its upper endpoint is either
$q+Lz\ge s$ or $L\ge s$.  The integrand is nonnegative, so this
enlargement can only increase the integral.

Rescale $h=Lv$.  With $\beta=b/L$ the integrand becomes
$k_\beta(v)\,dv$, with no additional factor of $L$.  The lower
endpoint $q_*/L$ and length $z$ satisfy
Lemma~\ref{thinning:interval} for this $x$ and
$\gamma=\sigma/p$.  Hence
\[
 \int_q^s\frac{(h-b)_+}{h^2}\,dh
 \le-\frac32\log(1-\gamma z)
 =-\frac32\log\left(1-\frac{d^2}{p^2}\right)
 =3\kappa.
\]
This includes $x=0$, which on the path can occur only at
$\lambda=1$, $x_c=0$.
\end{proof}

\begin{proof}[Proof of Theorem~\ref{thinning:payment}]
Combine the derivative reduction \eqref{thinning:derivative-loss}
with Lemma~\ref{thinning:window}:
\begin{align*}
 \lambda\Delta'-\Delta
 &\ge2\kappa-M\int_q^s\frac{(h-b)_+}{h^2}\,dh\\
 &\ge(2-3M)\kappa\ge0.
\end{align*}
Thus
$(\Delta/\lambda)'=(\lambda\Delta'-\Delta)/\lambda^2\ge0$.
At $\lambda=1$, the flag is deterministic, so $\nu=0$, $I=J_c$,
and $\Delta=F(J_c)$.  Comparing with this endpoint proves
\eqref{thinning:value}.  No limit as $\lambda$ tends to zero is needed.
\end{proof}

\subsection{The canonical high-entropy value}

\begin{theorem}[The canonical high-entropy value comparison]
\label{thinning:high-value}
Let $x\ge0$, $d>0$, and $x+d<1$.  Define
\[
 \bar h(x,d)=\frac{H(x+d)+H(x-d)}2,\qquad
 e_0(x,d)=(1-x)H\!\left(\frac d{1-x}\right).
\]
For every $e_0(x,d)\le e\le\bar h(x,d)$,
\begin{equation}
\label{thinning:high-value-bound}
 \mathcal K(x,d,e)\ge
 J_0(x,d,e):=F(H(x)-e)-F(\bar h(x,d)-e).
\end{equation}
\end{theorem}

\begin{proof}
The canonical classification at \eqref{canonical:transition} represents
every state in this closed high-entropy branch by
\eqref{thinning:path}, with a nonnegative core center and
$\mathcal K(x,d,e)=\lambda c(a_c,b_c)$.
Put $\alpha_c=1-x_c^2$.  Lemma~\ref{capacity:onebit} and convexity
of $F$, with $F(0)=0$, give
\[
 c(a_c,b_c)\ge\alpha_c F(J_c/\alpha_c)\ge F(J_c).
\]
Theorem~\ref{thinning:payment} therefore implies
\[
 J_0(x,d,e)=F(\nu+J)-F(\nu)
 \le\lambda F(J_c)
 \le\lambda c(a_c,b_c)=\mathcal K(x,d,e).
\]
\end{proof}

The canonical optimizer is used here only to evaluate a lower bound
on joining cost.  For an original pair law with means $x\pm d$ and
individual entropies $e_a,e_b$, the value of $e$ in
\eqref{thinning:high-value-bound} is their actual average.  The
original child information values remain $H(x+d)-e_a$ and
$H(x-d)-e_b$; their average is $\bar h(x,d)-e$.
Convexity of $F$ therefore gives the common floor at those original
child states:
\[
 \frac{F(H(x+d)-e_a)+F(H(x-d)-e_b)}2
 \ge F(\bar h(x,d)-e).
\]
No individual entropy moment of the canonical law is substituted for
either original child entropy.
